\midsectiontitle{Armas e Armaduras}

Em \textit{Overgrown}, cada arma é uma escolha de identidade.
Não existem regras para ditar quantas armas o aventureiro pode carregar ou quando utilizá-las, estar preparado
é parte do valor de sobreviver no Pináculo.

\begin{rpgboxtips}{Requisitos de Armas}
Cada arma possui um ou mais requisitos de atributos permanentes. Esses requisitos não dependem de DIV e podem ser cumpridos em qualquer nível. Bônus temporários, magias e estados não contam para esse fim.

Um personagem que não cumpra todos os requisitos ainda pode empunhar e atacar com a arma. Enquanto estiver empunhando ou utilizando essa arma, sofre \textbf{$-$10 em todos os testes que realizar}, mesmo quando o teste não estiver diretamente relacionado à arma. A penalidade termina assim que a arma deixa de ser empunhada ou utilizada. Rolagens de dano não são testes e não recebem essa penalidade.
\end{rpgboxtips}

\begin{rpgboxtips}{Configurações de Empunhadura Dupla}
Certos pares de armas formam uma \textbf{Configuração de Empunhadura Dupla}, apresentada em perfil próprio. Para usar esse perfil, o personagem deve empunhar uma cópia válida da arma em cada mão e cumprir o requisito da configuração.

Ao realizar um ataque, declare se está usando o perfil individual de uma das armas ou o perfil da configuração. O perfil de Empunhadura Dupla substitui o dano, a Ameaça e as Propriedades Especiais das armas individuais durante toda a ação. Benefícios dos dois perfis individuais não se acumulam com o perfil duplo. Ataques separados concedidos por uma configuração contam para o limite geral de ataques adicionais e não geram novos ataques extras. Uma configuração não pode usar Golpe Duplo, Disparo Duplo, Ataque Secundário ou outra fonte de ataque adicional na mesma ação.
\end{rpgboxtips}

% ==============================================================================
\titlesection{Armas Corpo a Corpo}
% ==============================================================================

\renewcommand{\tabularxcolumn}[1]{m{#1}}
\setlength{\extrarowheight}{3pt}
\renewcommand{\arraystretch}{1.2}
\begin{tabularx}{\textwidth}{|>{\columncolor{white}\centering\arraybackslash}m{0.18\textwidth}|>{\columncolor{white}\centering\arraybackslash}m{0.08\textwidth}|>{\columncolor{white}\centering\arraybackslash}m{0.06\textwidth}|>{\columncolor{white}\centering\arraybackslash}m{0.13\textwidth}|>{\columncolor{white}\centering\arraybackslash}m{0.12\textwidth}|>{\columncolor{white}\centering\arraybackslash}X|}
    \hline
    \multicolumn{1}{|c|}{\textbf{\ringm{Nome}}} &
    \multicolumn{1}{c|}{\textbf{\ringm{Ameaça}}} &
    \multicolumn{1}{c|}{\textbf{\ringm{Peso}}} &
    \multicolumn{1}{c|}{\textbf{\ringm{Dano Base}}} &
    \multicolumn{1}{c|}{\textbf{\ringm{Requisito}}} &
    \multicolumn{1}{c|}{\textbf{\ringm{Propriedade Especial}}} \\
    \hline
    Adaga & 18 & 1 & 1D6 + MOD de GRA & GRA 3 & \textbf{Assassin:} Contra um alvo Desprevenido, o primeiro acerto da cena aplica \link{status:ardiloso}{Sangramento Ardiloso}. \\
    \hline
    Alabarda & 19 & 5 & 2D8 + MOD de MIG & MIG 10 e SEN 5 & \textbf{Halberdier:} Alterne entre dano cortante e perfurante. Inimigos que avançarem a 2 unidades sofrem ataque de oportunidade. Acertos críticos aplicam \link{status:brutal}{Sangramento Brutal}. \\
    \hline
    Bo & 20 & 3 & 2D6 + MOD de MIG & MIG 3 e GRA 3 & \textbf{Monk:} Acertos críticos aplicam \link{status:concussionado}{Concussão}. Acertos comuns concedem +1 Esquiva até seu próximo turno (acumulável até +3). \\
    \hline
    Clava & 20 & 2 & 1D8 + MOD de MIG & MIG 2 & \textbf{Bruta:} Arma simples de impacto. Acertos críticos aplicam \link{status:concussionado}{Concussão} e \link{status:brutal}{Sangramento Brutal}. \\
    \hline
    Escudo & 20 & 3 & 1D6 + MOD de MIG & FOR 3 & \textbf{Guardian:} Ocupa uma mão e adiciona \textbf{+5 ao Valor de Bloqueio}. Apenas um escudo pode contribuir para cada Bloqueio. \\
    \hline
    Espada de Duas Mãos & 20 & 4 & 3D6 + MOD de MIG & MIG 10 e FOR 5 & \textbf{Guerreiro:} Adiciona \textbf{+3 ao Valor de Bloqueio}. Quando um alvo bloquear um ataque desta arma, reduza o Valor de Bloqueio total dele à metade, arredondando para baixo. \\
    \hline
    Espada de Uma Mão & 19 & 2 & 2D6 + MOD de MIG & MIG 5 & \textbf{Warden:} Use escudo sem penalidade de carga. \\
    \hline
    Lança & 19 & 3 & 2D8 + MOD de MIG & MIG 7 e SEN 3 & \textbf{Lancer:} Inimigos que avançarem a 1 unidade sofrem ataque de oportunidade. Pode ser arremessada (4 unidades), causando metade do dano. \\
    \hline
    Machado de Uma Mão & 19 & 3 & 2D6 + MOD de MIG & MIG 5 & \textbf{BloodLetter:} Acertos críticos aplicam \link{status:brutal}{Sangramento Brutal}. \\
    \hline
\end{tabularx}

\vspace{4pt}
\renewcommand{\tabularxcolumn}[1]{m{#1}}
\setlength{\extrarowheight}{3pt}
\renewcommand{\arraystretch}{1.2}
\begin{tabularx}{\textwidth}{|>{\columncolor{white}\centering\arraybackslash}m{0.18\textwidth}|>{\columncolor{white}\centering\arraybackslash}m{0.08\textwidth}|>{\columncolor{white}\centering\arraybackslash}m{0.06\textwidth}|>{\columncolor{white}\centering\arraybackslash}m{0.13\textwidth}|>{\columncolor{white}\centering\arraybackslash}m{0.12\textwidth}|>{\columncolor{white}\centering\arraybackslash}X|}
    \hline
    \multicolumn{1}{|c|}{\textbf{\ringm{Nome}}} &
    \multicolumn{1}{c|}{\textbf{\ringm{Ameaça}}} &
    \multicolumn{1}{c|}{\textbf{\ringm{Peso}}} &
    \multicolumn{1}{c|}{\textbf{\ringm{Dano Base}}} &
    \multicolumn{1}{c|}{\textbf{\ringm{Requisito}}} &
    \multicolumn{1}{c|}{\textbf{\ringm{Propriedade Especial}}} \\
    \hline
    Machado de Duas Mãos & 20 & 7 & 3D8 + MOD de MIG & MIG 12 e FOR 5 & \textbf{Berserker:} Todo acerto aplica \link{status:brutal}{Sangramento Brutal}. Com armadura leve ou sem armadura, causa 2D6 adicional. \\
    \hline
    Maça de Uma Mão & 19 & 3 & 2D6 + MOD de MIG & MIG 5 & \textbf{Crusader:} Com escudo, causa 1D8 adicional. Em acertos críticos, reduza em 2 o Valor de Bloqueio usado contra o ataque, até o mínimo de 0. \\
    \hline
    Maça de Duas Mãos & 20 & 6 & 2D8 + 1D6 + MOD de MIG & MIG 10 e FOR 5 & \textbf{Legionary:} Acertos críticos aplicam \link{status:atordoado}{Atordoado} por 1 turno. Ataques reduzem Resistência do alvo em 1 (máx. $-$5). \\
    \hline
    Martelo de Uma Mão & 19 & 3 & 2D6 + MOD de MIG & MIG 6 e FOR 3 & \textbf{Templar:} Com escudo, causa 1D8 adicional. Acertos aplicam \link{status:brutal}{Sangramento Brutal}. \\
    \hline
    Martelo de Duas Mãos & 20 & 6 & 2D10 + MOD de MIG & MIG 12 e FOR 5 & \textbf{Avenger:} Adiciona \textbf{+3 ao Valor de Bloqueio}. Acertos críticos aplicam \link{status:brutal}{Sangramento Brutal} e \link{status:concussionado}{Concussão} simultaneamente. \\
    \hline
    Rapieira & 18 & 2 & 2D6 + MOD de GRA & GRA 8 e SEN 3 & \textbf{Duelist:} Reduza em 2 PC o custo de Ataque Secundário com esta arma, até o mínimo de 3 PC. Se o alvo errar um ataque contra você, ganhe \textbf{Vantagem} no próximo ataque contra ele. \\
    \hline
    Soqueiras & 19 & 1 & 1D6 + MOD de GRA & GRA 8 e MIG 3 & \textbf{Pugilist:} Para cada 5 MOD de GRA, realize um ataque adicional ao usar uma Ação Padrão para atacar com as Soqueiras. Cada ataque adicional causa apenas 1D6, sem MOD ou dano adicional, não pode resultar em crítico, não ativa efeitos ao acertar e não gera outros ataques. Esta é uma exceção expressa ao limite geral de um ataque adicional. Ao avançar em direção a um inimigo para atacar, você \textbf{não provoca ataques de oportunidade} nesse deslocamento. \\
    \hline
    Tonfa & 19 & 2 & 1D8 + MOD de MIG & GRA 5 e FOR 5 & \textbf{Enforcer:} Quando seu Bloqueio reduzir ao menos 1 ponto de dano de um ataque corpo a corpo, receba +1 VB até o início do próximo turno. \\
    \hline
\end{tabularx}

% ==============================================================================
\titlesection{Armas de Longo Alcance}
% ==============================================================================

\vspace{4pt}
\renewcommand{\tabularxcolumn}[1]{m{#1}}
\setlength{\extrarowheight}{3pt}
\renewcommand{\arraystretch}{1.2}
\begin{tabularx}{\textwidth}{|>{\columncolor{white}\centering\arraybackslash}m{0.18\textwidth}|>{\columncolor{white}\centering\arraybackslash}m{0.08\textwidth}|>{\columncolor{white}\centering\arraybackslash}m{0.06\textwidth}|>{\columncolor{white}\centering\arraybackslash}m{0.13\textwidth}|>{\columncolor{white}\centering\arraybackslash}m{0.12\textwidth}|>{\columncolor{white}\centering\arraybackslash}X|}
    \hline
    \multicolumn{1}{|c|}{\textbf{\ringm{Nome}}} &
    \multicolumn{1}{c|}{\textbf{\ringm{Ameaça}}} &
    \multicolumn{1}{c|}{\textbf{\ringm{Peso}}} &
    \multicolumn{1}{c|}{\textbf{\ringm{Dano Base}}} &
    \multicolumn{1}{c|}{\textbf{\ringm{Requisito}}} &
    \multicolumn{1}{c|}{\textbf{\ringm{Propriedade Especial}}} \\
    \hline
    Arco Composto & 19 & 5 & 2D8 + MOD de GRA & GRA 8, MIG 5 e SEN 3 & \textbf{Hunter:} Ataques acima de 20 unidades têm Vantagem e +1D6. Pode canalizar magias como flechas. \\
    \hline
    Arco Leve & 20 & 1 & 1D8 + MOD de GRA & GRA 3 & \textbf{Scout:} Oculto ou primeiro a atacar, causa 1D8 extra. Pode ser empunhado e guardado como ação livre. \\
    \hline
    Arco Recurvo & 19 & 3 & 2D6 + MOD de GRA & GRA 6 e SEN 3 & \textbf{Ranger:} Em terreno elevado, causa 1D6 extra. Pode canalizar magias como flechas. \\
    \hline
    Besta Leve & 19 & 3 & 2D6 + MOD de GRA & GRA 3 e SEN 2 & \textbf{Pathfinder:} Seu mecanismo rúnico permite disparar sem gastar ação para recarregar. \\
    \hline
    Besta Pesada & 20 & 5 & 3D6 + MOD de GRA & GRA 6 e SEN 5 & \textbf{Marksman:} Marque um alvo no início do combate; ataques contra ele causam 1D8 extra. Exige ação de movimento para recarregar. \\
    \hline
    Bola de Cristal & 20 & 2 & 1D6 + MOD de WIS & WIS 8 e SEN 3 & \textbf{Channeler:} Adiciona metade do MOD de WIS (arredondado acima) aos testes de ataque mágico e o MOD de WIS completo ao dano de magias. \\
    \hline
    Bucaneira & 20 & 7 & 3D6 + MOD de MIG & MIG 10 e FOR 7 & \textbf{Buccaneer:} Acertos críticos ignoram o Valor de Bloqueio da armadura do alvo, mas não o bônus de escudos ou armas, e aplicam \link{status:brutal}{Sangramento Brutal}. Exige ação de movimento para recarregar. \\
    \hline
    Cajado & 20 & 3 & 2D6 + MOD de WIS & WIS 3 & \textbf{Arcanist:} Ao conjurar magias, gaste 10 IEP para aumentar o dano em 1D10. Pode ser usado em corpo a corpo sem penalidade. \\
    \hline
\end{tabularx}

\vspace{4pt}
\renewcommand{\tabularxcolumn}[1]{m{#1}}
\setlength{\extrarowheight}{3pt}
\renewcommand{\arraystretch}{1.2}
\begin{tabularx}{\textwidth}{|>{\columncolor{white}\centering\arraybackslash}m{0.18\textwidth}|>{\columncolor{white}\centering\arraybackslash}m{0.08\textwidth}|>{\columncolor{white}\centering\arraybackslash}m{0.06\textwidth}|>{\columncolor{white}\centering\arraybackslash}m{0.13\textwidth}|>{\columncolor{white}\centering\arraybackslash}m{0.12\textwidth}|>{\columncolor{white}\centering\arraybackslash}X|}
    \hline
    \multicolumn{1}{|c|}{\textbf{\ringm{Nome}}} &
    \multicolumn{1}{c|}{\textbf{\ringm{Ameaça}}} &
    \multicolumn{1}{c|}{\textbf{\ringm{Peso}}} &
    \multicolumn{1}{c|}{\textbf{\ringm{Dano Base}}} &
    \multicolumn{1}{c|}{\textbf{\ringm{Requisito}}} &
    \multicolumn{1}{c|}{\textbf{\ringm{Propriedade Especial}}} \\
    \hline
    Espingarda de Caça & 19 & 4 & 2D8 + MOD de GRA & GRA 7, MIG 5 e SEN 3 & \textbf{Trapper:} Use Mira Focada como ação livre, sem gastar PC. Até 5 unidades, causa 1D6 extra. Exige ação de movimento para recarregar. \\
    \hline
    Grimório & 19 & 2 & 1D8 + MOD de WIS & WIS 10 e SEN 5 & \textbf{Archmage:} Durante um Descanso Pleno, armazene no Grimório uma magia cuja fórmula esteja disponível ao personagem. Apenas uma magia pode ficar armazenada por vez; armazenar outra substitui a anterior. A magia armazenada pode ser conjurada mesmo fora das escolhas atuais do personagem, mas cada conjuração paga seu custo de IEP normalmente. Se estiver fora das escolhas atuais, seu custo de IEP é dobrado. O Grimório nunca elimina o custo da magia. \\
    \hline
    Pistola & 18 & 2 & 2D6 + MOD de GRA & GRA 5 e SEN 3 & \textbf{Gunslinger:} Recarregar a pistola exige uma ação de Movimento. \\
    \hline
    Rifle de Precisão & 19 & 8 & 2D10 + MOD de GRA & GRA 10 e SEN 10 & \textbf{Deadeye:} Acertos críticos causam 1D10 extra e ignoram o Valor de Bloqueio da armadura do alvo, mas não o bônus de escudos ou armas. Recarregar exige ação padrão completa. Desvantagem contra alvos a menos de 5 unidades. \\
    \hline
    Varinha & 19 & 1 & 1D8 + MOD de WIS & WIS 3 e SEN 2 & \textbf{Sorcerer:} Ataque básico causa Dano Mágico sem consumir IEP. Quando esse ataque causa ao menos 1 ponto de dano a um alvo vivo, aliado ou inimigo, recupere 1D4 IEP. Um mesmo alvo só pode alimentar esta recuperação uma vez por rodada, e a recuperação não pode elevar o IEP acima do máximo. Construtos, objetos, cadáveres e alvos sem vida não ativam o efeito. \\
    \hline
\end{tabularx}

% ==============================================================================
\titlesection{Configurações de Empunhadura Dupla}
% ==============================================================================

Os perfis abaixo são tratados como armas próprias durante a ação em que forem usados. O peso indicado é o peso total do par. Salvo indicação contrária, somente o primeiro ataque da configuração soma MOD ao dano.

\vspace{4pt}
\begin{tabularx}{\textwidth}{|>{\columncolor{white}\centering\arraybackslash}m{0.19\textwidth}|>{\columncolor{white}\centering\arraybackslash}m{0.08\textwidth}|>{\columncolor{white}\centering\arraybackslash}m{0.06\textwidth}|>{\columncolor{white}\centering\arraybackslash}m{0.14\textwidth}|>{\columncolor{white}\centering\arraybackslash}m{0.13\textwidth}|>{\columncolor{white}\centering\arraybackslash}X|}
    \hline
    \multicolumn{1}{|c|}{\textbf{\ringm{Configuração}}} &
    \multicolumn{1}{c|}{\textbf{\ringm{Ameaça}}} &
    \multicolumn{1}{c|}{\textbf{\ringm{Peso}}} &
    \multicolumn{1}{c|}{\textbf{\ringm{Dano Base}}} &
    \multicolumn{1}{c|}{\textbf{\ringm{Requisito}}} &
    \multicolumn{1}{c|}{\textbf{\ringm{Propriedade Especial}}} \\
    \hline
    Adagas Gêmeas & 19 & 2 & Dois ataques de 1D6 & GRA 8 & \textbf{Twin Assassin:} Os dois ataques sofrem apenas $-$1 de acerto. Se ambos acertarem o mesmo alvo, aplique uma carga de \link{status:ardiloso}{Sangramento Ardiloso}. Aplique no máximo uma carga por ação. \\
    \hline
    Espadas Gêmeas & 19 & 4 & Dois ataques de 1D8 & MIG 7 e GRA 5 & \textbf{Twin Warden:} Os dois ataques sofrem $-$2 de acerto. Se ambos acertarem o mesmo alvo, cause +1D6 de dano após resolver o segundo ataque. \\
    \hline
    Machados Gêmeos & 20 & 6 & Dois ataques de 1D8 & MIG 8 e GRA 5 & \textbf{Twin BloodLetter:} Os dois ataques sofrem $-$2 de acerto. Se ambos causarem dano ao mesmo alvo, o segundo aplica \link{status:brutal}{Sangramento Brutal}. \\
    \hline
    Tonfas Gêmeas & 19 & 4 & 2D6 + MOD de MIG & GRA 7 e FOR 5 & \textbf{Twin Enforcer:} Adicione +3 VB. Uma vez por rodada, quando um Bloqueio reduzir dano de ataque corpo a corpo, gaste 2 PC para contra-atacar imediatamente. Esse ataque não exige nova Reação e não gera ataques adicionais. \\
    \hline
    Bestas Leves Gêmeas & 20 & 6 & Dois ataques de 1D8 & GRA 7 e SEN 5 & \textbf{Twin Pathfinder:} Os ataques sofrem $-$2 de acerto e podem escolher alvos diferentes. Os mecanismos rúnicos independentes não exigem ação de recarga. \\
    \hline
    Pistolas Gêmeas & 19 & 4 & Dois ataques de 1D8 & GRA 8 e SEN 5 & \textbf{Twin Gunslinger:} Os ataques sofrem $-$1 de acerto e devem escolher o mesmo alvo. Uma única Ação de Movimento recarrega ambas as pistolas. \\
    \hline
\end{tabularx}

% ==============================================================================
\titlesection{Armaduras e Proteções}
% ==============================================================================

Em \textit{Overgrown}, armaduras não alteram a Resistência passiva. Elas fornecem um \textbf{Valor de Bloqueio (VB)}, usado para reduzir o dano de ataques que já acertaram. Armaduras mais pesadas exigem treinamento físico permanente, enquanto a Armadura de Couro depende de mobilidade e controle corporal.

\begin{itemize}
    \item \textbf{Valor de Bloqueio (VB):} quantidade fixa de dano que a armadura reduz quando o personagem declara um Bloqueio.
    \item \textbf{Requisito:} valor mínimo permanente de atributo necessário para usar a armadura corretamente. Bônus temporários não permitem cumprir requisitos.
    \item \textbf{Requisito não cumprido:} a armadura ainda ocupa sua Carga e aplica suas penalidades, mas não concede VB nem efeitos positivos e não pode ser usada para Bloquear.
\end{itemize}

\vspace{4pt}
\renewcommand{\tabularxcolumn}[1]{m{#1}}
\setlength{\extrarowheight}{3pt}
\renewcommand{\arraystretch}{1.2}
\begin{tabularx}{\textwidth}{|>{\columncolor{white}\centering\arraybackslash}m{0.20\textwidth}|>{\columncolor{white}\centering\arraybackslash}m{0.06\textwidth}|>{\columncolor{white}\centering\arraybackslash}m{0.08\textwidth}|>{\columncolor{white}\centering\arraybackslash}m{0.12\textwidth}|>{\columncolor{white}\centering\arraybackslash}X|}
    \hline
    \multicolumn{1}{|c|}{\textbf{\ringm{Nome}}} &
    \multicolumn{1}{c|}{\textbf{\ringm{Peso}}} &
    \multicolumn{1}{c|}{\textbf{\ringm{VB}}} &
    \multicolumn{1}{c|}{\textbf{\ringm{Requisito}}} &
    \multicolumn{1}{c|}{\textbf{\ringm{Efeitos}}} \\
    \hline
    Vestes de Pano      & 2 & 0 & Nenhum & \textbf{Condutor Rúnico:} +1 Espaço de Melhoria Rúnica. \textbf{Liberdade:} Sem limite de Esquiva. Sem penalidade de ruído. \\
    \hline
    Gambeson            & 2 & 3 & Nenhum & \textbf{Amortecida:} Reduz o dano de \link{status:concussionado}{Concussão} em 1D4. Sem penalidade de Esquiva ou ruído. \\
    \hline
    Armadura de Couro   & 4 & 5 & GRA 10 & \textbf{Agilidade:} +5 de Esquiva adicional. \textbf{Reforço:} Proteção básica sem prejuízo à mobilidade. \\
    \hline
    Cota de Malha       & 7 & 8 & MIG 12 & \textbf{Malha Trançada:} +2 RD contra dano cortante. \textbf{Barulhento:} $-$3 em testes de Furtividade. \\
    \hline
    Couraça de Metal    & 10 & 12 & MIG 15 & \textbf{Peso Morto:} $-$7 de Esquiva. \textbf{Impermeável:} Imune a \link{status:concussionado}{Concussão}. \\
    \hline
    Armadura de Placas  & 15 & 15 & MIG 20 & \textbf{Tanque de Guerra:} $-$15 de Esquiva. \\
    \hline
\end{tabularx}

\vspace{8pt}
\begin{rpgboxtips}{O Uso do Bloqueio}
Depois que um ataque perceptível acertar você, mas antes da rolagem de dano, gaste uma \textbf{Reação} para declarar um \textbf{Bloqueio}. Some o VB da sua armadura aos bônus fixos de um escudo, arma ou habilidade. O resultado reduz o dano daquele ataque, até o mínimo de 0.

\[\text{Dano final}=\max(0,\text{dano após modificadores}-\text{VB}-\text{RD})\]

Um Bloqueio reduz apenas uma instância de dano. Uma habilidade com vários dados e um único teste de acerto é uma instância; ataques com testes de acerto separados exigem Bloqueios separados. O Bloqueio e a Esquiva não podem ser usados contra a mesma instância.
\end{rpgboxtips}

\begin{rpgboxtips}{Sequência e Limites do Bloqueio}
Críticos acertam normalmente, mas podem ter seu dano reduzido. Reduzir o dano a 0 não transforma o acerto em erro: efeitos que ocorram \textbf{ao acertar} ainda são aplicados, mas efeitos que exijam \textbf{causar dano} não são ativados. VB excedente não é armazenado nem transferido para outro ataque.

Podem ser bloqueados ataques corpo a corpo, projéteis e efeitos de área perceptíveis que causem dano direto. Dano contínuo, ambiental, mental ou Puro não pode ser bloqueado, salvo indicação expressa. Contra uma área, cada personagem usa sua própria Reação e reduz apenas o dano que receberia; não existe teste de Bloqueio.
\end{rpgboxtips}

\begin{rpgboxtips}{Armaduras e Escudos}
Apenas uma armadura equipada fornece VB e efeitos. Um Gambeson pode ser vestido sob Cota de Malha, Couraça de Metal ou Armadura de Placas: nesse caso, seu peso e o efeito Amortecida permanecem, e seu VB é somado. Nenhuma outra combinação de armaduras acumula benefícios.

Um escudo ocupa uma mão e fornece \textbf{+5 VB} ao Bloqueio, inclusive quando o usuário veste Pano ou nenhuma armadura. O bônus de apenas um escudo pode ser aplicado a cada Bloqueio. Empunhar ou guardar um escudo segue as mesmas regras de uma arma de uma mão.
\end{rpgboxtips}
